\section{Full Evaluation and Reporting Protocol}
\label{app:evaluation}

\subsection{Required immutable run identity}

Every result row must resolve to a run record containing repository commit,
complete composed configuration, source and split manifest digests, seed,
checkpoint, training log, final metrics, and prediction artifact.  A group
record binds the three or more seed runs used for a reported mean and standard
deviation.  Failed/non-finite runs remain in the group; they may not be removed
after inspecting results.

\subsection{Joint-objective result template}

\begin{table}[h]
  \centering
  \caption{Mandatory joint-objective comparison.  Values are intentionally
  omitted until immutable three-seed run records exist; consistency is a
  diagnostic and cannot determine adoption by itself.}
  \label{tab:result_template}
  \scriptsize
  \begin{tabular}{@{}lrrrrrrrr@{}}
    \toprule
    Variant & KP $\downarrow$ & Pose-RP $\downarrow$ & $t$ m $\downarrow$ &
    $R$ deg $\downarrow$ & $f$ rel. $\downarrow$ & Line $\uparrow$ &
    Seg $\uparrow$ & Cons. $\downarrow$ \\
    \midrule
    direct-all & --- & --- & --- & --- & --- & --- & --- & --- \\
    joint-both & --- & --- & --- & --- & --- & --- & --- & --- \\
    stopgrad-pose & --- & --- & --- & --- & --- & --- & --- & --- \\
    stopgrad-dense & --- & --- & --- & --- & --- & --- & --- & --- \\
    \bottomrule
  \end{tabular}
\end{table}



In addition to this compact main table, the run bundle reports mean/median and
per-channel keypoint distance, translation axes, rotation, focal relative and
log error, pose-only reprojection, line Dice, segmentation mIoU, dense--pose
consistency, invalid-depth rate, eligible-point counts, parameter count, step
time, memory, and non-finite events.  Synthetic and real-image metrics remain
separate.

\subsection{Alignment and dataset manifest report}

For every reconstructed scene, a dataset paper table must be generated from the
published manifest rather than copied from configuration defaults.  Required
fields are accepted/rejected render counts; trajectory groups by split; camera
height, radius, azimuth, field-of-view, and adjacent-step distributions;
target-court counts; coverage/supervision-eligible-point histograms; render failures; and
the exact alignment artifact digest.  Rejected alignment attempts are reported
as negative evidence and contribute zero dataset samples.

\subsection{Failure cases}

Qualitative results must include at least: weak or occluded line evidence,
partial court views, neighbouring-court confusion, mid-plane near/far
rejection, predicted negative depth, a low consistency/high GT-error example,
3DGS extrapolation artifacts, and a real image outside the synthetic camera
envelope.  Cherry-picked successful overlays alone do not test the claimed
global mechanism.
